% Part: methods
% Chapter: proofs
% Section: example-1

\documentclass[../../../include/open-logic-section]{subfiles}

\begin{document}

\olfileid[hi]{mth}{prf}{ex1}

\olsection{एक उदाहरण}

हमारा पहला उदाहरण समुच्चयों के सम्मिलनों और सर्वनिष्ठों के बारे में
निम्नलिखित सरल तथ्य है। इससे परिभाषाओं को खोलना, संयोजनों और सार्वत्रिक
दावों के प्रमाण तथा स्थितियों द्वारा प्रमाण स्पष्ट होंगे।

\begin{prop}
किन्हीं समुच्चयों $A$, $B$ और $C$ के लिए, $A \cup (B \cap C) = (A \cup
B) \cap (A \cup C)$।
\end{prop}

आइए इसे सिद्ध करें!{}

\begin{proof}
हम दिखाना चाहते हैं कि किन्हीं समुच्चयों $A$, $B$ और $C$ के लिए, $A \cup
(B \cap C) = (A \cup B) \cap (A \cup C)$।
\begin{quote}
पहले हम प्रतिज्ञप्ति के कथन में ``$=$'' की परिभाषा खोलते हैं। याद करें कि
समुच्चयों को समान सिद्ध करने का अर्थ यह दिखाना है कि उनके !!{element} समान
हैं। अर्थात् $A \cup (B \cap C)$ के सभी !!{element}, $(A \cup B) \cap
(A \cup C)$ के भी !!{element} हैं, और इसके विपरीत। ``इसके विपरीत'' का
अर्थ है कि $(A \cup B) \cap (A \cup C)$ का प्रत्येक !!{element}, $A
\cup (B \cap C)$ का भी !!a{element} होना चाहिए। इसलिए परिभाषा खोलने पर
हम देखते हैं कि हमें एक संयोजन सिद्ध करना है। इसे लिख लें:
\end{quote}
परिभाषा से, $A \cup (B \cap C) = (A \cup B) \cap (A \cup C)$ तभी और
केवल तभी जब $A \cup (B \cap C)$ का प्रत्येक !!{element}, $(A \cup B)
\cap (A \cup C)$ का भी !!a{element} हो, और $(A \cup B) \cap (A \cup
C)$ का प्रत्येक !!{element}, $A \cup (B \cap C)$ का !!a{element} हो।
\begin{quote}
चूँकि यह एक संयोजन है, हमें प्रत्येक संयोज्य अलग-अलग सिद्ध करना होगा। पहले
से आरंभ करें: सिद्ध करें कि $A \cup (B \cap C)$ का प्रत्येक !!{element},
$(A \cup B) \cap (A \cup C)$ का भी !!a{element} है।

यह एक सार्वत्रिक दावा है, इसलिए हम $A \cup (B \cap C)$ के किसी मनमाने
!!{element} पर विचार करके दिखाते हैं कि वह $(A \cup B) \cap (A \cup
C)$ का भी !!a{element} अवश्य है। इस मनमाने !!{element} को संबोधित करने
के लिए कोई चर चुनेंगे, मान लें~$z$। हमारा प्रमाण आगे बढ़ता है:
\end{quote}
पहले हम सिद्ध करते हैं कि $A \cup (B \cap C)$ का प्रत्येक !!{element},
$(A \cup B) \cap (A \cup C)$ का भी !!a{element} है। मान लें $z \in A
\cup (B \cap C)$। हमें दिखाना है कि $z \in (A \cup B) \cap (A \cup
C)$।
\begin{quote}  
अब $\cup$ और~$\cap$ की परिभाषाएँ खोलने का समय है। उदाहरण के लिए, $\cup$
की परिभाषा है: $A \cup B = \Setabs{z}{z \in A \text{ अथवा } z \in B}$।
जब हम परिभाषा को ``$A \cup (B \cap C)$'' पर लागू करते हैं, तो परिभाषा
के ``$B$'' की भूमिका अब ``$B \cap C$'' निभाता है; इसलिए $A \cup (B
\cap C) = \Setabs{z}{z \in A \text{ अथवा } z \in B \cap C}$। अतः हमारी
मान्यता $z \in A \cup (B \cap C)$ का अर्थ है: $z \in \Setabs{z}{z \in
A \text{ अथवा } z \in B \cap C}$। और $z \in \Setabs{z}{\dots z\dots}$
तभी और केवल तभी जब \dots $z$ \dots, अर्थात् इस मामले में $z \in A$ या
$z \in B \cap C$।
\end{quote}
परिभाषा $\cup$ से, या तो $z \in A$ या $z \in B \cap C$।
\begin{quote}
चूँकि यह एक वियोजन है, स्थितियों द्वारा प्रमाण लागू करना उपयोगी होगा। हम
दोनों स्थितियाँ लेते हैं और दिखाते हैं कि प्रत्येक में हमारा अभीष्ट निष्कर्ष
(अर्थात् ``$z \in (A \cup B) \cap (A \cup C)$'') प्राप्त होता है।
\end{quote}
स्थिति 1: मान लें कि $z \in A$।
\begin{quote}
हमारी मान्यताओं के आधार पर आगे काम करने को अधिक कुछ नहीं है। इसलिए देखें कि
निष्कर्ष में क्या उपलब्ध है। हम दिखाना चाहते हैं कि $z \in (A \cup B)
\cap (A \cup C)$। $\cap$ की परिभाषा के आधार पर इसे दिखाने के लिए हमें
दिखाना होगा कि $z$, $(A \cup B)$ और $(A \cup C)$ दोनों में है। किंतु $z
\in A \cup B$ तभी और केवल तभी जब $z \in A$ या $z \in B$, और हमारे पास
पहले ही (स्थिति~1 की मान्यता के रूप में) $z \in A$ है। इसी तर्क से---$B$
के स्थान पर $C$ रखकर---$z \in A \cup C$। यह तर्क उलटी दिशा में चला था,
अतः अपने प्रमाण के लिए आवश्यक दिशा में इसे लिखें।
\end{quote}
चूँकि $z \in A$, इसलिए $z \in A$ या $z \in B$, और अतः परिभाषा~$\cup$
से $z \in A \cup B$। इसी प्रकार $z \in A \cup C$। किंतु परिभाषा~$\cap$
से इसका अर्थ है कि $z \in (A \cup B) \cap (A \cup C)$।
\begin{quote}
इससे स्थितियों द्वारा प्रमाण की पहली स्थिति पूरी होती है। अब हम दूसरी स्थिति
में निष्कर्ष निकालना चाहते हैं, जहाँ $z \in B \cap C$।
\end{quote}
स्थिति 2: मान लें कि $z \in B \cap C$।
\begin{quote}
फिर हम दो समुच्चयों के सर्वनिष्ठ के साथ काम कर रहे हैं। परिभाषा~$\cap$ लागू
करें:
\end{quote}
चूँकि $z \in B \cap C$, परिभाषा~$\cap$ से $z$, $B$ और $C$ दोनों का
!!a{element} अवश्य है।
\begin{quote}
अब अपने निष्कर्ष को फिर देखें। हमें दिखाना है कि $z$, $(A \cup B)$ और
$(A \cup C)$ दोनों में है। फिर समाधान तत्काल है।
\end{quote}
चूँकि $z \in B$, इसलिए $z \in (A \cup B)$। चूँकि $z \in C$, इसलिए $z
\in (A \cup C)$ भी। अतः $z \in (A \cup B) \cap (A \cup C)$।
\begin{quote}
यहाँ हमने $\cup$ और $\cap$ की परिभाषाएँ फिर लागू कीं, किंतु चूँकि हम उन
परिभाषाओं को पहले ही याद कर चुके हैं और पहले ही दिखा चुके हैं कि यदि $z$ दो
समुच्चयों में से एक में हो तो वह उनके सम्मिलन में है, इसलिए हमें अपने किए को
उतने विस्तार से बताने की आवश्यकता नहीं।

हमने स्थितियों द्वारा प्रमाण की दूसरी स्थिति पूरी कर ली है, इसलिए अब अपना
पहला निष्कर्ष कह सकते हैं।
\end{quote}
अतः यदि $z \in A \cup (B \cap C)$, तो $z \in (A \cup B) \cap (A \cup
C)$।
\begin{quote}
अब हमें केवल दूसरी दिशा दिखानी है कि $(A \cup B) \cap (A \cup C)$ का
प्रत्येक !!{element}, $A \cup (B \cap C)$ का !!a{element} है। पहले की
तरह हम प्रथम समुच्चय का कोई मनमाना !!{element} अपने पास मानकर और यह दिखाकर
कि वह दूसरे समुच्चय में अवश्य है, इस सार्वत्रिक दावे को सिद्ध करते हैं। हम
जो करने वाले हैं उसे लिखें।
\end{quote}
अब मान लें कि $z \in (A \cup B) \cap (A \cup C)$। हम दिखाना चाहते हैं
कि $z \in A \cup (B \cap C)$।
\begin{quote}
अब हम परिकल्पना $z \in (A \cup B) \cap (A \cup C)$ से काम कर रहे हैं।
आशा है कि प्रमाण के पहले भाग वाला ही~$z$ यहाँ प्रयोग करना अधिक उलझन पैदा
नहीं करता। उस भाग को पूरा करते ही वहाँ की हमारी सभी मान्यताएँ प्रभावहीन हो
गईं, इसलिए अब हम $z$ के बारे में नई मान्यताएँ बना सकते हैं। यदि यह उलझाता
हो तो आगे $z$ के स्थान पर कोई अन्य चर रख दें।

परिभाषा~$\cap$ से हम जानते हैं कि $z$, $A \cup B$ और $A \cup C$ दोनों
में है। और $\cup$ की परिभाषा से इसे आगे खोल सकते हैं: या तो $z \in A$ या
$z \in B$, और साथ ही या तो $z \in A$ या $z \in C$। यह फिर स्थितियों
द्वारा प्रमाण जैसा दिखता है---सिवाय इसके कि ``और'' इसे उलझा देता है। आप सोच
सकते हैं कि इसका अर्थ तीन संभावनाएँ होना है: $z$, $A$, $B$ या $C$ में है।
लेकिन यह भूल होगी। हमें सावधान रहना है, अतः प्रत्येक वियोजन पर बारी-बारी से
विचार करें।
\end{quote}
परिभाषा~$\cap$ से $z \in A \cup B$ और $z \in A \cup C$। परिभाषा
$\cup$ से $z \in A$ या $z \in B$। हम स्थितियों में भेद करते हैं।
\begin{quote}
चूँकि हमारा ध्यान पहले वियोजन पर है, हमने अभी अपना दूसरा वियोजन ($A \cup
C$ को खोलने से) प्राप्त नहीं किया। वास्तव में अभी उसकी आवश्यकता नहीं। पहली
स्थिति $z \in A$ है, और किसी समुच्चय का !!a{element}, उस समुच्चय का किसी
अन्य समुच्चय के साथ सम्मिलन का भी !!a{element} होता है। अतः स्थिति~1 सरल
है:
\end{quote}
स्थिति 1: मान लें कि $z \in A$। इससे $z \in A \cup (B \cap C)$ निष्पन्न
होता है।
\begin{quote}
अब दूसरी स्थिति $z \in B$ लें। यहाँ हम दूसरा $\cup$ खोलेंगे और एक और बार
स्थितियों द्वारा प्रमाण करेंगे:
\end{quote}
स्थिति 2: मान लें कि $z \in B$। चूँकि $z \in A \cup C$, इसलिए या तो $z
\in A$ या $z \in C$। हम आगे स्थितियों में भेद करते हैं:

स्थिति 2a: $z \in A$। तब फिर $z \in A \cup (B \cap C)$।
\begin{quote}
ठीक है, यह कुछ विचित्र था। इस स्थिति के लिए वास्तव में हमें मान्यता~$z \in
B$ की आवश्यकता नहीं थी, किंतु कोई समस्या नहीं।
\end{quote}
स्थिति 2b: $z \in C$। तब $z \in B$ और $z \in C$, अतः $z \in B \cap
C$, और फलतः $z \in A \cup (B \cap C)$।
\begin{quote}
इससे स्थितियों द्वारा दोनों प्रमाण पूरे होते हैं, और इसलिए दूसरा आधा भी
पूरा हुआ।
\end{quote}
अतः यदि $z \in (A \cup B) \cap (A \cup C)$, तो $z \in A \cup (B \cap
C)$।
\end{proof}

\end{document}
